\lecture{8.4}{Longitudinal IPTW Formula}
\label{lec:8-4}

\subsection{Generalized Propensity Scores}

For a point treatment, the propensity score is the conditional probability of
receiving the observed or target treatment. In a longitudinal study, there is a
time-specific treatment mechanism at every decision time.

\begin{definition}[Generalized propensity score]
For a binary treatment and fixed regime $\abar=(a_0,\ldots,a_K)$, define
\[
  g_k(a_k\given \lbar_k,\abar_{k-1})
  =
  \Prob(A_k=a_k\given \Lbar_k=\lbar_k,\Abar_{k-1}=\abar_{k-1}).
\]
For the always-treated regime, write
\[
  g_k(\lbar_k)
  =
  \Prob(A_k=1\given \Lbar_k=\lbar_k,\Abar_{k-1}=\mathbf{1}_k).
\]
\end{definition}

The time-specific propensity score is local: it concerns the treatment decision
at time $k$ among individuals whose previous treatment history is compatible
with the regime.

\subsection{Composite Treatment Probabilities}

The probability of receiving an entire treatment profile is the product of the
time-specific probabilities. For regime $\abar$ define
\[
  \gbar_K(\Lbar_K)
  =
  \prod_{k=0}^K
  g_k(a_k\given \Lbar_k,\Abar_{k-1}=\abar_{k-1}).
\]
For the always-treated regime this is
\[
  \gbar_K(\Lbar_K)
  =
  g_0(L_0)g_1(\Lbar_1)\cdots g_K(\Lbar_K).
\]

This product is the probability that a person with the realized covariate
history would follow the regime through all treatment times, under the observed
treatment mechanism.

\subsection{Treatment-Regime Weights}

The inverse probability weight for a regime $\abar$ is
\[
  H_{\abar}
  =
  \frac{\ind(\Abar_K=\abar_K)}{\gbar_K(\Lbar_K)}.
\]
For the always-treated regime,
\[
  H_{\mathbf{1}}
  =
  \frac{A_0A_1\cdots A_K}
       {g_0(L_0)g_1(\Lbar_1)\cdots g_K(\Lbar_K)}.
\]

The numerator selects individuals whose observed treatment history matches the
regime. The denominator upweights those individuals by the inverse probability
of observing that full history.

\begin{figure}[H]
\centering
\begin{tikzpicture}[
  node distance=.9cm,
  box/.style={draw, rounded corners=2pt, minimum width=2.35cm, minimum height=.72cm},
  arr/.style={-{Latex[length=2.2mm]}, thick}
]
\node[box] (p0) {$g_0$};
\node[box, right=of p0] (p1) {$g_1$};
\node[right=of p1] (dots) {$\cdots$};
\node[box, right=of dots] (pK) {$g_K$};
\node[box, below=of dots, minimum width=6.4cm] (prod) {$\gbar_K=\prod_{k=0}^K g_k$};
\node[box, below=of prod, minimum width=6.4cm] (w) {$H_{\abar}=\ind(\Abar=\abar)/\gbar_K$};
\draw[arr] (p0) -- (prod);
\draw[arr] (p1) -- (prod);
\draw[arr] (pK) -- (prod);
\draw[arr] (prod) -- (w);
\end{tikzpicture}
\caption{Longitudinal IPTW uses the product of time-specific treatment probabilities.}
\label{fig:ch8-longitudinal-iptw}
\end{figure}

\subsection{IPTW Identification for Treatment Profiles}

\begin{identificationbox}
Under consistency, sequential exchangeability, and longitudinal positivity,
\[
  \E\{Y(\abar)\}
  =
  \E\left[
    \frac{\ind(\Abar_K=\abar_K)}
         {\prod_{k=0}^K g_k(a_k\given \Lbar_k,\Abar_{k-1}=\abar_{k-1})}
    Y
  \right].
\]
For the always-treated regime,
\[
  \E\{Y(1,\ldots,1)\}
  =
  \E\left[
    \frac{A_0A_1\cdots A_K}
         {g_0(L_0)g_1(\Lbar_1)\cdots g_K(\Lbar_K)}
    Y
  \right].
\]
\end{identificationbox}

The formula creates a pseudo-population in which following the regime is no
longer associated with measured history. Individuals with rare regime-compatible
histories receive larger weights because they represent many similar individuals
who did not follow the regime.

\subsection{Positivity Problems Over Time}

Longitudinal weights can become unstable quickly. If $K=5$ and each
time-specific probability equals $0.5$, then the composite probability is
$0.5^6=0.015625$. The corresponding weight is $64$ for an individual who
follows the whole regime. If some probabilities are smaller, weights can be much
larger.

There are two kinds of positivity problems:
\begin{itemize}
  \item \textbf{Structural violations}: a treatment sequence is impossible for
  some history, such as a contraindicated drug after severe toxicity.
  \item \textbf{Practical violations}: the sequence is possible in principle but
  rarely observed in the sample.
\end{itemize}
Both problems are more common as the number of treatment times grows.

\subsection{Equivalence with G-Computation}

The IPTW and g-computation formulas identify the same counterfactual mean. The
equivalence follows by repeated use of iterated expectation.

For the always-treated regime, let
\[
  \Qbar_{K+1}(\Lbar_K)=\E(Y\given \Abar_K=\mathbf{1}_{K+1},\Lbar_K).
\]
Then
\[
\begin{aligned}
&\E\left[
  \frac{A_0\cdots A_K}{g_0(\Lbar_0)\cdots g_K(\Lbar_K)}Y
\right] \\
&=
\E\left[
  \frac{A_0\cdots A_K}{g_0(\Lbar_0)\cdots g_K(\Lbar_K)}
  \E(Y\given \Lbar_K,\Abar_K)
\right] \\
&=
\E\left[
  \frac{A_0\cdots A_K}{g_0(\Lbar_0)\cdots g_K(\Lbar_K)}
  \Qbar_{K+1}(\Lbar_K)
\right].
\end{aligned}
\]
Now condition on $\Lbar_K,\Abar_{K-1}$. Since
\[
  \E(A_K\given \Lbar_K,\Abar_{K-1}=\mathbf{1}_K)=g_K(\Lbar_K),
\]
the factor $A_K/g_K(\Lbar_K)$ averages to one among those compatible with the
past regime. Therefore the last expression becomes
\[
  \E\left[
  \frac{A_0\cdots A_{K-1}}{g_0(\Lbar_0)\cdots g_{K-1}(\Lbar_{K-1})}
  \E\{\Qbar_{K+1}(\Lbar_K)\given \Lbar_{K-1},\Abar_{K-1}=\mathbf{1}_K\}
  \right].
\]
The conditional expectation is $\Qbar_K(\Lbar_{K-1})$. Repeating the same
argument removes $A_{K-1}/g_{K-1}$, then $A_{K-2}/g_{K-2}$, and so on, yielding
\[
  \E\{\Qbar_1(L_0)\},
\]
which is the g-computation formula.

\begin{keyidea}
IPTW and g-computation are two representations of the same identifying
functional. G-computation integrates outcome regressions forward through the
covariate process; IPTW reweights observed regime followers to represent the
target population.
\end{keyidea}
